% Class-compatibility preamble retained for a reader excerpt of
%   Johan Hastad, Bjorn Martinsson, Tamio-Vesa Nakajima, Stanislav Zivny,
%   "A logarithmic approximation of linearly ordered colourings",
%   arXiv:2404.19556v6, CC BY 4.0.
%
% The source is a plain article; the only compatibility work here is to repeat the
% author's own declaration block mechanically and unchanged (pinned-original-source.tex
% lines 5 and 18-23: the ams packages, the plain/definition theorem styles and the
% author's four theorem environments), and to carry the two packages that the publisher
% layout and the author's cross-references need (hyperref and cleveref, the latter with
% the author's own nameinlink option). The author's font package newtx is deliberately
% not reproduced: it is a font choice for the journal layout and the excerpt is rendered
% in the delivery baseline fonts, which changes no mathematical content.
% No mathematical macro, statement, symbol or assumption is introduced, changed or
% removed by this file.
\documentclass[10pt,a4paper]{article}
\usepackage{amssymb,amsmath,amsthm}
\usepackage{hyperref}
\usepackage[nameinlink]{cleveref}

\theoremstyle{plain}
\newtheorem{theorem}{Theorem}
\newtheorem{lemma}[theorem]{Lemma}
\newtheorem{observation}[theorem]{Observation}
\theoremstyle{definition}
\newtheorem{remark}[theorem]{Remark}
\title{A randomized-and-derandomized subprocedure that finds a large vertex set meeting every edge of a 2-3-uniform hypergraph in the required parity}
\author{Johan H\{\textbackslash{}aa\}stad, Bj\textbackslash{}"orn Martinsson, Tamio-Vesa Nakajima, Stanislav \textbackslash{}v\{Z\}ivn\textbackslash{}'y}
\date{arxiv:2404.19556v6; CC BY 4.0}
\begin{document}
\maketitle
\noindent\textit{Editorial note.} All delivered mathematical text is the author's own, copied from the pinned author source sources/combinatorics-fl05-a-logarithmic-approximation-of-lin/pinned-original-source.tex; no statement or proof step was rewritten, reordered or removed, and the author's own spelling is preserved verbatim.\par\medskip
\noindent\textit{Editorial note.} Preamble provenance: this card uses sources/combinatorics-fl05-a-logarithmic-approximation-of-lin/retained-preamble.tex, which repeats the author's ams packages, the plain and definition theorem styles and the author's four newtheorem declarations (pinned source lines 5 and 18-23) mechanically and unchanged, and adds only hyperref and cleveref (with the author's own nameinlink option) because the author's cross-reference commands are used by other parts of the paper. The author's font package newtx is deliberately not reproduced: it is a layout font choice, the excerpt is rendered in the delivery baseline fonts, and no mathematical content changes.\par\medskip
\noindent\textit{Editorial note.} Bibliography provenance: the source has no thebibliography environment; it ends with bibliographystyle alphaurl plus bibliography bib (author source lines 354-355), and the paper names its bibliography file bib which is not staged anywhere in this repository, so this card deliberately counts only a statement whose statement and proof cite nothing. This card carries no bibliography and no entry was reconstructed or invented from memory.\par\medskip
\noindent\textit{Editorial note.} The counted statement and its proof contain no citation and no cross-reference, so the card needs neither a bibliography excerpt nor a navigation stub.\par\medskip
\noindent\textit{Allowed standard prerequisites (editorial).}\begin{itemize}
\item The linearly ordered (LO) colouring of a hypergraph and the class of LO k-colourable hypergraphs. Defined in the introduction of the paper, before source line 100; the counted statement uses the case k = 2. Fixes the notion of an LO 2-colouring that the input hypergraph is assumed to admit.
\item The linear system Av = 1\textasciicircum{}m attached to an LO 2-colourable hypergraph, and the statement that it does not fix the value of any variable. Set up in the section introducing the algorithm based on equations modulo 2 (source lines 147-150); the hypothesis of the counted statement is exactly this non-fixing condition. Supplies the affine solution space in which the randomised choice of a\_i is made, and the fact that for every vertex x some homogeneous solution v\textasciicircum{}i has v\textasciicircum{}i\_x = 1.
\end{itemize}\medskip
\begin{lemma}\label{lem:innerstep}
    There is an algorithm which, if given a $\{2, 3\}$-uniform hypergraph $H$ with $n$ vertices and
    $m$ edges that admits an LO 2-colouring and such that the implied linear
    system of equations $Av=1^m$ does not fix the value of any variable, outputs a subset $T$ of vertices that intersects edges of size three in zero or two vertices and edges of size two in exactly one vertex.
    Moreover, we have  $\vert T\vert \geq n / 2$. The algorithm runs in $O(n^3 + nm)$ time.
\end{lemma}

\begin{proof}

    We first describe a randomised version of our algorithm, and then derandomise
    it. The set of solutions to $Av=1^m$ is an affine space and hence a generic solution can be written as $v=v^0+\sum_{i=1}^r a_i v^i$ for a basic solution $v^0$, linearly independent solutions to the homogeneous system $v^i$, and field elements (in this case bits) $a_i$.
    The fact that no variable is fixed implies that for each vertex $x$ there is 
    some positive $i$ such that $v^i_x=1$.
    
     For the randomised algorithm choose $a_1, \ldots, a_r$ to be independent identically distributed uniformly random bits, and set $T$ to be 
     the set of variables $x$, such that $v_x=0$.  Clearly $T$ satisfies
     the conclusion of the lemma as in each edge we have an odd number of ones. 
     For every vertex $x$, there exists positive $i$ such that $v_x^i = 1$ --- hence, due to the influence of $a_i v_x^i$, we see that $x$ is included in $T$ with probability $\frac{1}{2}$.
     Thus, on average $T$ contains half the vertices.

    Now, we derandomise this algorithm using the method of conditional
    expectations. Go through the variables $a_i$ in increasing order and 
    fix its value once and for all.  Fixing the value of $a_i$ determines
    the value(s) of some $v_x$ while other values remain undetermined.  For each
    value being determined $v_x^i=1$ and hence one value of $a_i$ gives
    the final value 0 and the other gives final value 1.  Set $a_i$ such
    that at least half the determined values are 0.  After we have 
    fixed all $a_i$ this way, we have a final solution with at least
    $n/2$ zeroes.
    
    The bottleneck of the running time of this algorithm is solving the linear system of equations. This can be done in the advertised running time since every equation has $O(1)$ entries. 
\end{proof}

\end{document}
